% arxiv-style preprint, single-column, compiles with pdflatex.
% \documentclass{article} + an inline preamble so it builds without
% external arxiv.sty.

\documentclass[10pt]{article}

% --- packages --------------------------------------------------------------
\usepackage[utf8]{inputenc}
\usepackage[T1]{fontenc}
\usepackage{lmodern}
\usepackage[a4paper, margin=1in]{geometry}
\usepackage{microtype}
\usepackage{hyperref}
\hypersetup{colorlinks=true, linkcolor=black, citecolor=black, urlcolor=black}
\usepackage{url}
\usepackage{booktabs}
\usepackage{amsmath}
\usepackage{amssymb}
\usepackage{amsfonts}
\usepackage{nicefrac}
\usepackage{graphicx}
\usepackage{xcolor}
\usepackage{titlesec}
\usepackage[numbers,square,sort&compress]{natbib}
\usepackage{enumitem}
\usepackage{caption}
\usepackage{xspace}
\usepackage{parskip}

\setlength{\parskip}{6pt}
\setlength{\parindent}{0pt}

% --- arxiv-preprint title block --------------------------------------------
\titleformat*{\section}{\Large\bfseries}
\titleformat*{\subsection}{\large\bfseries}
\titleformat*{\subsubsection}{\normalsize\bfseries}

\title{\Large\bfseries When Procedural-Skill SFT Works:\\
  Negative-Result-Driven Iteration Across 0.8B--4B Qwen3.5 Models}

\author{Italo Strozzi\thanks{\texttt{istrozzi@matematica.ufrj.br}}\\
  Independent Researcher}

\date{May 6, 2026}

% --- shorthand --------------------------------------------------------------
\newcommand{\BL}{\textsc{baseline}\xspace}
\newcommand{\CU}{\textsc{curated}\xspace}
\newcommand{\Delt}{\ensuremath{\Delta}\xspace}
\newcommand{\Qweno}{Qwen3.5-0.8B\xspace}
\newcommand{\Qwent}{Qwen3.5-2B\xspace}
\newcommand{\Qwenf}{Qwen3.5-4B\xspace}

% ----------------------------------------------------------------------------
\begin{document}
\maketitle

\begin{abstract}
We study whether supervised fine-tuning (SFT) on a small corpus of (task + injected procedural-skill block, Opus chain-of-thought, judge-confirmed pass) demonstrations teaches a smaller language model to apply procedural skills \emph{differentially} when shown an explicit procedure. Across eight training-recipe variants and three model sizes (\Qweno, \Qwent, \Qwenf) on a 200-task / 40-skill holdout, we find: (i) at 0.8B the SFT signal is dominated by format learning and never produces a clean \CU-vs-\BL{} lift across five recipe iterations; (ii) at 2B, the same simplest recipe (LoRA $r{=}16$, 353 rows, three epochs) yields \CU{} pass-rate $0.825$ vs \BL{} $0.750$, and a controlled attribution split (pre-SFT 2B at $0.685/0.710$) isolates $+11.5$\,pp \CU{} lift to SFT contribution and $+14.5$\,pp to base-model scaling; (iii) at 4B, the same recipe lifts \CU{} to $0.880$ with $\Delta$ shrinking to $+0.045$ (still individually significant, $p{=}0.049$), although the v1.9$\to$v2.0 $\Delta$-shrinkage is consistent with but not statistically supported as bench saturation at this sample size. We additionally find that the bench scoring is format-sensitive in a way that prevents direct measurement of pre-SFT 4B reasoning, that the v1.7 ``mode collapse'' failure mode is a 0.8B capacity-floor artifact rather than a structural property of the SFT shape, and that the calibration regression visible in 2B SFT generality probes is eliminated at 4B. We release pipeline, recipes, eval scripts, and full episode-level logs.
\end{abstract}

\section{Introduction}

A common compositional-skills hypothesis holds that LLM capabilities can be decomposed into reusable skills, and that explicit procedural instructions can elicit better task performance from smaller students than implicit-skill priming alone (cf.\ Skill-Mix \citep{yu2024skillmix}, procedural-injection benchmarks). A natural follow-up: does demonstrating the explicit procedure during fine-tuning teach a smaller student to \emph{use} such procedures more effectively at inference time, or merely to imitate the response shape?

We answer this empirically. We construct a self-contained pipeline that generates a procedural-skill SFT corpus from a hand-curated declarative skill catalog, train Qwen3.5 students at three scales with eight recipe variants, evaluate against a 200-task holdout with deterministic and LLM-judge-based scoring, and probe out-of-distribution behavior with a 52-prompt generality battery.

The main empirical results are:
\begin{enumerate}[leftmargin=*]
  \item \textbf{At 0.8B, SFT teaches format, not procedural-skill differentiation.} Five recipe variants (varying corpus composition, chat-template patches, partial fine-tuning, skill-block stripping) all cluster post-SFT \CU{} pass-rate inside a 2\,pp band of $0.565$--$0.585$, statistically indistinguishable from \BL.
  \item \textbf{At 2B, the simplest recipe works.} \CU{} pass-rate $0.825$ vs \BL{} $0.750$. With a pre-SFT 2B control ($0.685/0.710$), the \CU{} lift decomposes cleanly into $+14.5$\,pp from base-model scaling and $+11.5$\,pp from SFT contribution. The procedural-$\Delta$ ($\CU - \BL$) lifts from $+0.025$ (pre-SFT) to $+0.075$ (post-SFT).
  \item \textbf{At 4B, the same recipe yields a smaller but still significant $\Delta$.} \CU{} $0.880$, \BL{} $0.835$, $\Delta$ $+0.045$ (McNemar exact $p{=}0.049$). The shrinkage from v1.9's $+0.075$ is directionally consistent with bench saturation but not statistically distinguishable from noise at $n{=}200$ ($p{=}0.225$ one-sided bootstrap; § \ref{sec:stats}). The pre-SFT 4B attribution control fails because the deterministic judge requires literal \texttt{ANSWER:} lines that the base model does not reliably emit, so the 4B SFT contribution is bounded but not directly measurable on this bench.
  \item \textbf{Mode-collapse failure modes are capacity-floor artifacts.} A v1.7 partial-FT recipe at 0.8B produced an apparent $+0.150$ $\Delta$ but with collapsed \BL{} ($0.465$); the same SFT shape at 4B produces no analogous failure ($0/52$ phantom skill-block references in general chat).
  \item \textbf{Generality is preserved.} A 52-prompt OOD probe shows v2.0 (4B+LoRA) matches base 4B on factual recall ($34/34$ HIT), eliminates a calibration regression observed in v1.9 (2B+LoRA), and produces no format lock-in despite three epochs of training on a single response shape.
\end{enumerate}

The negative-result iteration sequence (five 0.8B variants before model-size pivot) and the format-sensitivity finding (the bench is partially a format-compliance test for non-SFT-trained students) are themselves contributions, since they constrain the interpretation of similar SFT-on-procedural-demonstrations claims.

\section{Pipeline and Bench}

\subsection{Skill catalog and procedural rewrite}
We hand-curate 40 declarative skills across categories (logical fallacies, deductive forms, cognitive biases, interpretive patterns, spatial reasoning), seeded from Skill-Mix Table 5/6 \citep{yu2024skillmix} and expanded across three review batches. Each declarative entry contains \texttt{\{name, definition, example\}}. We use Opus 4.7 (1M context) to rewrite each into a procedural form: ordered procedure steps, when-to-use criteria, constraints, and a positive/negative worked example. The resulting \texttt{bench\_catalog/skills.json} contains the procedural skill definitions used as injection content.

\subsection{Task synthesis}
A task synthesizer (we call it \textsc{s3.m5}) prompts Opus, given a procedural skill, to produce 15 tasks designed to require that skill. Tasks carry one of four query types: \texttt{YES\_NO} (122/200, 61\%), \texttt{FREE\_FORM} (33/200, 16.5\%), \texttt{SINGLE\_WORD} (23/200, 11.5\%), \texttt{RANKING} (22/200, 11\%). Tasks are split into \texttt{tasks\_train.json} (400 tasks) and \texttt{tasks\_eval.json} (200 tasks), disjoint by task-uid and ordered by a \texttt{task\_skill\_map} that records the 1:1 task-to-skill pairing.

\subsection{SFT corpus generation}
We trace each training task with Opus under the \textsc{solver} system prompt with the corresponding procedural skill block injected. Episodes that pass the LLM-judge are filtered to a curated-only SFT corpus of 353 rows. Each row is a chat-format triple: \texttt{(system: SOLVER + skill block, user: task, assistant: Opus's chain-of-thought + ANSWER: line)}.

\subsection{Evaluation protocol}
For every \texttt{(task, condition)} pair on the 200-task holdout we generate one student response and judge it. Two conditions: \BL{} (\textsc{solver} prompt without skill block) and \CU{} (\textsc{solver} prompt with the corresponding procedural skill injected). Judge dispatch: deterministic \texttt{ANSWER:}-line extraction for \texttt{YES\_NO}/\texttt{SINGLE\_WORD}/\texttt{RANKING}, Opus 4.7 LLM-judge for \texttt{FREE\_FORM}. Decoding: greedy ($T{=}0$), \texttt{repetition\_penalty}\,$=1.05$ (to break degenerate token loops on out-of-distribution prompts), \texttt{max\_new\_tokens} 1024 (post-SFT) or 2048 (pre-SFT base; thinking-mode is verbose). The two reported per-condition statistics are pass-rate (binary judge verdict) and mean score (0--1 judge score). We also report $\Delta = \mathrm{pass\_rate}(\CU) - \mathrm{pass\_rate}(\BL)$, the differential procedural lift.

\subsection{Training}
LoRA via TRL\,$\geq$\,0.18 \citep{trl} and PEFT\,$\geq$\,0.13. Adapters target all-linear modules. The chat template is patched with \texttt{\{\% generation \%\}\dots\{\% endgeneration \%\}} markers to enable assistant-only loss masking (TRL $\geq$\,0.18 falls back to loss-on-all-tokens otherwise). For partial-FT runs (v1.7, v1.8) we freeze all but the top-$N$ transformer layers + \texttt{lm\_head} + final \texttt{norm}, and use bitsandbytes 8-bit Adam \citep{dettmers2023qlora} to fit \Qweno{} on consumer hardware.

\section{Recipe Iterations}
We report eight recipe variants. Five at 0.8B (v1--v1.8) explore corpus composition, partial fine-tuning, and skill-block masking; two at 2B/4B (v1.9, v2.0) test model-size as a separate axis.

\begin{table}[h]
\centering
\small
\begin{tabular}{@{}llllll@{}}
\toprule
\textbf{Variant} & \textbf{Base} & \textbf{Recipe} & \textbf{BL} & \textbf{CU} & \textbf{$\Delta$} \\
\midrule
Pre-SFT 0.8B & qwen3.5:0.8b (Ollama) & none & 0.510 & 0.565 & $+0.055$ \\
Pre-SFT 2B   & Qwen/Qwen3.5-2B (HF)  & none & 0.685 & 0.710 & $+0.025$ \\
Pre-SFT 4B   & Qwen/Qwen3.5-4B (HF)  & none & $0.000$\,$\dagger$ & $0.000$\,$\dagger$ & --- \\
Pre-SFT haiku-4-5 (ref) & claude-haiku-4-5 & none & 0.785 & 0.800 & $+0.015$ \\
\addlinespace
\textbf{v1}    & 0.8B & LoRA $r{=}8$, curated-only, 353 rows           & 0.635 & 0.585 & $-0.050$ \\
\textbf{v1.5}  & 0.8B & LoRA + BL+CU mixed corpus                       & 0.650 & 0.425 & $-0.225$ \\
\textbf{v1.6}  & 0.8B & LoRA + drop ceiling skills + chat-patch         & 0.645 & 0.565 & $-0.080$ \\
\textbf{v1.7}  & 0.8B & partial FT (top-6 layers + lm\_head)            & 0.465 & 0.615 & $+0.150$\,$\ddagger$ \\
\textbf{v1.8}  & 0.8B & partial FT + skill-block stripped from training & 0.545 & 0.570 & $+0.025$ \\
\textbf{v1.9}  & 2B   & LoRA $r{=}16$, curated-only, 353 rows           & \textbf{0.750} & \textbf{0.825} & $\boldsymbol{+0.075}$ \\
\textbf{v2.0}  & 4B   & LoRA $r{=}32$, curated-only, 353 rows           & \textbf{0.835} & \textbf{0.880} & $\boldsymbol{+0.045}$ \\
\bottomrule
\end{tabular}
\caption{Recipe-by-recipe pass rates on the 200-task holdout, two conditions per task. $\dagger$ Pre-SFT 4B's $0/0$ is a format-compliance artifact (\S\ref{sec:format-sensitivity}). $\ddagger$ v1.7's apparent $+\Delta$ is an artifact of \BL{} collapse, not differential learning (\S\ref{sec:v17-mode-collapse}).}
\label{tab:headline}
\end{table}

\section{0.8B: Five Negative Results}

\subsection{v1: simplest recipe, format-only learning}
LoRA $r{=}8$/$\alpha{=}16$ on the 353-row \CU{} corpus for three epochs at $\mathrm{lr}{=}2{\times}10^{-4}$ lifts the 0.8B \BL{} pass-rate by $+12.5$\,pp ($0.510 \to 0.635$) but does not lift \CU; $\Delta$ flips from $+0.055$ (pre-SFT) to $-0.050$ (post-SFT). Per-skill inspection at $n{=}200$ shows the lift is concentrated on tasks where the model previously failed to produce the bench's \texttt{ANSWER:} format and now does so. We interpret v1 as having taught primarily the response shape (step-by-step + \texttt{ANSWER:} line), with no detectable transfer of procedural skill use.

\subsection{v1.5--v1.6: corpus composition does not unlock differential lift}
Adding baseline-passing rows to the SFT corpus (v1.5; teaching the model to produce procedural responses both with and without skill-block presence) collapses \CU{} pass-rate to $0.425$. Dropping ceiling-passing skills (v1.6; removing tasks where the base already passes $\geq\!0.90$ to avoid the model learning to add procedure overhead on tasks it can solve directly) lands at $0.645/0.565$ --- equivalent to v1 within noise.

\subsection{v1.7: full-rank partial FT and the mode-collapse artifact}\label{sec:v17-mode-collapse}
Replacing LoRA with full-rank fine-tuning of the top-6 layers + \texttt{lm\_head} + final \texttt{norm} (using 8-bit Adam to fit GPU budget) produces an apparent positive $\Delta$ ($+0.150$) but with \BL{} pass-rate \emph{collapsed} to $0.465$ --- substantially below pre-SFT $0.510$. Inspection of v1.7 \BL{} failures reveals the model emits ``\texttt{Step 1 (SKILL block: \dots)}'' prefaces and procedural step language even when no skill block is in the system prompt --- i.e.\ the model has learned to expect a skill block and hallucinates one when absent. The $+0.150$ $\Delta$ does not reflect the model learning to use skill blocks; it reflects the model losing the ability to function without one.

\subsection{v1.8: skill-block stripping fixes mode collapse, no new ceiling}
Stripping the skill block from the system prompt of every training row (so the model never sees a skill block during training, only the procedural response) recovers \BL{} ($0.545$) but lands \CU{} at $0.570$ --- within the $0.565$--$0.585$ cluster shared by v1, v1.6, and v1.8. The mode-collapse failure mode is fixed; no new ceiling is unlocked.

\subsection{0.8B summary}
Three well-formed variants (v1, v1.6, v1.8) cluster \CU{} within 2\,pp. The bench ceiling at 0.8B does not move with corpus composition, recipe, chat-template patches, or partial fine-tuning. We infer the ceiling is set by base-model capacity, not by the training recipe.

\section{2B and 4B: Model Size as the Decisive Axis}

\subsection{v1.9 (2B): SFT works}
Same recipe shape as v1 (LoRA, curated-only, 353 rows, 3 epochs), bumped to $r{=}16$/$\alpha{=}32$ to give the larger base more LoRA capacity, on \Qwent. Result: \BL{} $0.750$ / \CU{} $0.825$ / $\Delta$ $+0.075$. \CU{} pass-rate exceeds the pre-SFT haiku-4-5 reference of $0.800$.

A pre-SFT 2B control ($0.685/0.710$) cleanly attributes the lift over pre-SFT 0.8B (Table~\ref{tab:attribution}). The procedural-$\Delta$ lifts from $+0.025$ (pre-SFT) to $+0.075$ (post-SFT), a $+5$\,pp differential procedural-application lift attributable to SFT.

This is the cleanest available measurement of the original SFT hypothesis. SFT does teach differential procedural use; the 0.8B failures were base-capacity-bound.

\begin{table}[h]
\centering
\small
\begin{tabular}{@{}lrr@{}}
\toprule
& $\Delta$BL & $\Delta$CU \\
\midrule
Base scaling 0.8B $\to$ 2B  & $+0.175$ & $+0.145$ \\
SFT contribution at 2B      & $+0.065$ & $\mathbf{+0.115}$ \\
\midrule
Total v1.9 vs pre-SFT 0.8B  & $+0.240$ & $+0.260$ \\
\bottomrule
\end{tabular}
\caption{v1.9 attribution split. Base scaling is measured between the two pre-SFT controls; SFT contribution is the residual from pre-SFT 2B to v1.9.}
\label{tab:attribution}
\end{table}

\subsection{v2.0 (4B): bench saturation and the format-sensitivity artifact}\label{sec:format-sensitivity}
Same recipe at LoRA $r{=}32$/$\alpha{=}64$ on \Qwenf. Result: \BL{} $0.835$ / \CU{} $0.880$ / $\Delta$ $+0.045$. The \CU{} $> 0.85$ absolute gate is cleared; the $\Delta \geq 5$\,pp gate is a statistical tie with v1.9 within $n{=}200$ binomial noise ($\sim 3$\,pp per condition, $\sim 5$\,pp on $\Delta$).

The pre-SFT 4B attribution control was attempted but failed: base \Qwenf{} reasons correctly on bench tasks (spot-checks confirm step-by-step reasoning that reaches the right answer) but does not reliably emit the literal \texttt{ANSWER:} format the deterministic judge requires. With 88.5\% of the bench using the deterministic judge, pre-SFT 4B scores $\sim\!0/200$, which is a format-compliance artifact rather than a reasoning measurement. The v2.0 SFT contribution at 4B is therefore bounded by the saturation trend (plausibly $0.03$--$0.10$ \CU) but not directly measured on this bench.

The v1.9$\to$v2.0 deltas ($\Delta$BL $+0.085$, $\Delta$CU $+0.055$, $\Delta\Delta -0.030$) are unaffected by this artifact since both v1.9 and v2.0 are SFT-trained format-compliant students. The shrinkage of the SFT-attributable $\Delta$ from $+0.075$ to $+0.045$ is consistent with bench saturation; § \ref{sec:stats} shows it is not statistically distinguishable at $n{=}200$.

\subsection{Statistical tests of the per-model and cross-model $\Delta$ estimates}\label{sec:stats}

Each task is run in both \BL{} and \CU{} conditions, so the verdicts are paired. We run McNemar's test (continuity-corrected $\chi^2$ for $n_{\mathrm{discordant}} \geq 25$, exact binomial otherwise) per model, and a paired bootstrap (10{,}000 resamples of \texttt{task\_uid} with replacement) for 95\% CIs on $\Delta$. Cross-model: we resample v1.9 and v2.0 independently and report a one-sided bootstrap test for $H_0\!: \Delta_{\mathrm{v1.9}} \leq \Delta_{\mathrm{v2.0}}$ (the direction predicted by the saturation hypothesis).

\begin{table}[h]
\centering
\small
\begin{tabular}{@{}lccc@{}}
\toprule
& $\Delta$ & McNemar two-sided $p$ & bootstrap 95\% CI on $\Delta$ \\
\midrule
v1.9 & $+0.0750$ & $\chi^2{=}5.60,\ p{=}0.018^{*}$  & $[+0.020,\ +0.130]$ \\
v2.0 & $+0.0450$ & exact, $p{=}0.049^{*}$            & $[+0.005,\ +0.085]$ \\
\bottomrule
\end{tabular}
\caption{Per-model paired tests on the procedural-$\Delta$. Both models have a positive $\Delta$ that is individually distinguishable from zero at $\alpha{=}0.05$ ($^{*}$), although v2.0's $p{=}0.049$ would not survive Bonferroni correction across the two tests.}
\label{tab:stats-per-model}
\end{table}

For the cross-model saturation test:
\begin{itemize}[leftmargin=*,nosep]
  \item $\Delta_{\mathrm{v1.9}} - \Delta_{\mathrm{v2.0}} = +0.030$
  \item Paired-bootstrap 95\% CI on the difference: $[-0.040,\ +0.100]$
  \item One-sided bootstrap $p$ for $H_0\!: \Delta_{\mathrm{v1.9}} \leq \Delta_{\mathrm{v2.0}}\!: p = 0.225$
\end{itemize}

The CI on the difference includes zero and the one-sided $p$ is $0.225$. The $\Delta$-shrinkage from v1.9 to v2.0 is consistent with bench saturation in direction but is \emph{not statistically distinguishable from sampling noise at the present sample size}. Confirming saturation would require either (a) larger $n$ ($\sim 800$--$1200$ per condition for $80\%$ power to detect an effect of this magnitude), (b) a bench where the procedural lift is 2--3$\times$ larger so the absolute $\Delta$ is more robustly measurable, or (c) multi-seed evaluation to reduce per-condition variance.

\subsection{Generality probe: v2.0 preserves out-of-distribution behavior}
We probe v1.9, v2.0, base 2B, and base 4B with a 52-prompt OOD battery (10 generality + 42 factual including 8 false-premise ``trick'' prompts) under a generic system prompt (not the \textsc{solver} prompt; Table~\ref{tab:generality}).

\begin{table}[h]
\centering
\small
\begin{tabular}{@{}lcccc@{}}
\toprule
& format-locked & factual HIT & trick CONFAB (eyeballed) & phantom skill block \\
\midrule
v1.9 (2B) & $0/52$ & $33/34$\,$\dagger$ & $6/8$ & $0/52$ \\
base 4B   & $0/52$ & $34/34$ & $\sim\!1/8$ & $0/52$ \\
v2.0 (4B) & $0/52$ & $34/34$ & $\mathbf{0/8}$ & $0/52$ \\
\bottomrule
\end{tabular}
\caption{Generality probe under generic system prompt. $\dagger$ v1.9 swapped \emph{Brave New World} attribution from Aldous Huxley to H. G. Wells --- a plausible-author confabulation absent from base 4B and v2.0.}
\label{tab:generality}
\end{table}

v1.9 introduced a small calibration regression on adversarial prompts: the SFT-induced ``produce a confident step-by-step answer'' prior leaks into out-of-distribution behavior, making v1.9 commit to plausible-sounding inventions where base 2B would refuse. v2.0 does not exhibit this --- the 4B base has enough native capacity to recognize false premises and refuse cleanly without the SFT shape overriding refusal.

The bench-time finding that 17.5\% of v2.0's \BL{} responses reference a ``skill block'' not present in the system prompt does \emph{not} extend to general-chat distributions: under a generic system prompt, the SFT-induced skill-block reflex is dormant ($0/52$). The phantom-mention is contextual to SFT-distribution-similar prompts.

\section{Discussion}

\subsection{What did SFT actually teach?}
At 2B, the cleanest measurement we have, SFT contributed:
\begin{itemize}[leftmargin=*,nosep]
  \item $+11.5$\,pp absolute \CU{} pass-rate, on top of base 2B's already-respectable $0.710$
  \item $+5$\,pp differential procedural-application lift ($\Delta$ from $+0.025$ to $+0.075$), evidence that the model learned to \emph{use} the procedure rather than just imitate the response shape
  \item $+6.5$\,pp absolute \BL{} pass-rate, attributable to format learning (\texttt{ANSWER:}-line emission, step-by-step structure)
\end{itemize}
The SFT contribution decomposes into roughly half ``format'' and half ``procedure'' by these crude metrics. The two are not independent --- producing the procedural-step shape implicitly enforces a checking discipline that may itself be the procedural lift --- but the differential $\Delta$-lift confirms there is a procedure-vs-no-procedure component of the gain that pure format-learning cannot explain.

\subsection{Bench format-sensitivity is itself a methodological finding}
The pre-SFT 4B attribution control failure surfaces a property of the bench scoring that is invisible when only SFT-trained students are evaluated: 88.5\% of the holdout uses a deterministic \texttt{ANSWER:}-line extractor, which heavily penalizes models that reason correctly but emit free-form conclusions. Base \Qwenf{} is such a model. The bench's $0/200$ for pre-SFT 4B should be read as ``format-compliance baseline $\approx 0$'', not ``reasoning $\approx 0$''.

This has interpretive consequences for any paper claiming ``$+X$\,pp from SFT on procedural demonstrations'' without a base-model control benched under format-tolerant scoring. The post-SFT lift conflates format learning, format-conditional reasoning gating, and procedural application; isolating these requires either (a) a stricter base-model evaluation prompt that hits the format target, (b) a format-tolerant judge for base controls, or (c) \texttt{FREE\_FORM}-only evaluation. We did not pursue these in the present study, but mark them as required steps for future replication.

\subsection{Mode collapse reinterpreted}
The v1.7 ``mode collapse'' --- partial-FT producing apparent $+\Delta$ at the cost of \BL{} integrity --- appeared at 0.8B as a structural concern about SFT-on-procedural-demonstrations and motivated the v1.8 skill-block stripping intervention. Spot-checks of v2.0 \BL{} responses show the same phantom skill-block reference ($35/200$, 17.5\%) but with $94.3\%$ of those responses still passing --- at 4B, the phantom mention is cosmetic rather than load-bearing. The model has the capacity to execute the procedure regardless of whether it remembers seeing a skill block in the system prompt.

This re-frames mode collapse as a base-capacity-floor artifact: at 0.8B, the model's SFT-absorbed expectation of seeing a skill block crashed performance when one was absent; at 4B, the same expectation produces a confabulated mention but does not derail the answer. The v1.8 skill-block-stripping fix is unnecessary at the model sizes where SFT actually delivers.

\subsection{Negative results were informative}
The five 0.8B variants are individually publishable as negative results --- corpus composition, partial fine-tuning, and chat-template patching all failed to lift the SFT signal at this scale. They are jointly informative: together they constrain the diagnosis to ``model size is the binding constraint,'' which the 2B run validates. We argue that running through this negative-result series before moving to a larger model was epistemically efficient ($\sim$\$15 in compute, three days of GPU time) compared to skipping straight to 4B with no priors on what fails.

\section{Limitations}

\begin{enumerate}[leftmargin=*]
  \item \textbf{Pre-SFT 0.8B was run via Ollama}, not the same HuggingFace transformers path used for 2B/4B. Quantization and chat-template differences may inflate or deflate its baseline relative to a strictly-comparable HF run. We did not re-bench 0.8B under the same path.
  \item \textbf{Pre-SFT 4B is unmeasured on this bench.} The format-compliance artifact prevents direct attribution of v2.0's lift between base scaling and SFT contribution.
  \item \textbf{$n{=}200$ per condition limits the resolution of the $\Delta$ statistic.} Per-condition binomial standard error is $\sim\!3$\,pp; per-model McNemar's tests confirm v1.9 ($p{=}0.018$) and v2.0 ($p{=}0.049$) each have a positive $\Delta$ distinguishable from zero, but the cross-model bootstrap test for the saturation hypothesis fails ($p{=}0.225$ one-sided, 95\% CI on $\Delta_{\mathrm{v1.9}} - \Delta_{\mathrm{v2.0}}$ is $[-0.040, +0.100]$). The $\Delta$ shrinkage is directional but not statistically supported at this $n$.
  \item \textbf{The bench tests single-skill-per-task application, not skill composition.} The original procedural-skills hypothesis posits compositional benefit (apply skill $A$ then skill $B$); the present bench cannot distinguish ``model uses procedure when shown'' from ``model would have applied roughly the right pattern anyway, because the task was synthesized to require exactly that skill.''
  \item \textbf{The SFT corpus is small (353 rows).} Whether scaling the corpus to $10\times$ would change the saturation curve at 4B is untested.
  \item \textbf{Attribution at 4B is bounded but not measured.} We rely on saturation-trend extrapolation for the 4B SFT contribution range ($0.03$--$0.10$ \CU).
  \item \textbf{The judge is an LLM (Opus 4.7).} While we patched the judge for deterministic-extraction robustness and re-judged $\sim\!17\%$ of episodes that hit transient auth failures, we have not cross-validated against a second judge or human raters. Judge bias toward the SFT-induced response shape is plausible but unmeasured.
\end{enumerate}

\section{Conclusion}

Procedural-skill SFT works as a recipe for teaching a smaller language model to apply explicit procedures more effectively when shown --- but only at base sizes where the model already produces format-compliant reasoning. At 0.8B, five recipe variants establish that the SFT signal cannot exceed format-learning at this capacity. At 2B, the simplest recipe (LoRA $r{=}16$, 353 rows, 3 epochs) produces a measurable $+11.5$\,pp \CU{} lift attributable to SFT contribution and a $+5$\,pp differential procedural-$\Delta$ lift, with clean attribution from a pre-SFT 2B control. At 4B, the same recipe lifts to \CU{} $0.880$ with $\Delta$ at $+0.045$ (still individually distinguishable from zero, McNemar $p{=}0.049$); the directional shrinkage from v1.9's $+0.075$ is consistent with bench saturation but does not survive a paired-bootstrap test at $n{=}200$ ($p{=}0.225$ one-sided), so the saturation claim is suggestive rather than confirmed and motivates either larger-$n$ runs or harder benches for follow-up work.

The negative-result iteration sequence and the bench format-sensitivity finding both constrain the interpretation of similar SFT-on-procedural-demonstrations claims. We argue future work in this area should (i) include a pre-SFT base-model control evaluated under format-tolerant scoring, (ii) measure compositional skill application rather than single-skill template-matching, and (iii) bound the SFT contribution within statistical noise at the chosen sample size before claiming model-size-conditional results.

\section*{Acknowledgments}
Compute provided by vast.ai (\Qwent{} and \Qwenf{} training and partial evaluation) and a local consumer GPU (RTX 3090\,Ti, evaluation and probing). Skill catalog content seeded from Skill-Mix \citep{yu2024skillmix} Tables 5 and 6 and expanded by hand. Opus 4.7 (1M context) was used for skill procedural rewriting, task synthesis, baseline trace generation, and judge-of-record scoring.

\section*{Code and data availability}
Repository: \url{https://github.com/izzortsi/skillmix-n-skillsbench} (branch \texttt{dev}). Episode-level eval logs at \texttt{data/pipeline-runs/default/bench-eval-*}. Engineering reports in \texttt{b1.reports/} (cited in Appendix~\ref{app:reports}).

% --- bibliography ----------------------------------------------------------
\begin{thebibliography}{9}

\bibitem{yu2024skillmix}
D.~Yu, S.~Kaur, A.~Gupta, J.~Brown-Cohen, A.~Goyal, S.~Arora.
\textit{Skill-Mix: A Flexible and Expandable Family of Evaluations for AI Models}.
NeurIPS 2024.

\bibitem{arora2023theory}
S.~Arora, A.~Goyal.
\textit{A Theory for Emergence of Complex Skills in Language Models}.
2023.

\bibitem{li2026skillsbench}
Li et al.
\textit{SkillsBench: Benchmarking the Effectiveness of Skill Injection on LLMs}.
2026.

\bibitem{hu2022lora}
E.~J.~Hu, Y.~Shen, P.~Wallis, Z.~Allen-Zhu, Y.~Li, S.~Wang, L.~Wang, W.~Chen.
\textit{LoRA: Low-Rank Adaptation of Large Language Models}.
ICLR 2022.

\bibitem{dettmers2023qlora}
T.~Dettmers, A.~Pagnoni, A.~Holtzman, L.~Zettlemoyer.
\textit{QLoRA: Efficient Finetuning of Quantized LLMs}.
NeurIPS 2023.

\bibitem{trl}
L.~Tunstall, E.~Beeching, et al.
\textit{TRL: Transformer Reinforcement Learning library}, version $\geq$\,0.18.

\end{thebibliography}

\appendix

\section{Pipeline file structure}\label{app:pipeline}
\begin{verbatim}
skillmix-n-skillsbench/
|-- core/                                  shared schema + provider abstractions
|-- harness/                               LinearAgent: streaming + thinking-block capture
|-- b2_benchmarks/skillsbench/             corpus_harness, llm_judge, procedural_catalog
|-- s1_extracting_skill_names/             declarative-skill seeders (Wikipedia-style)
|-- s2_extracting_skills_from_text/        procedural-rewrite stage
|-- s3_generating_skill_examples/m5_*.py   task synthesis from procedural skills
|-- s4_extracting_skill_usage_instances/   solver + trace + judge composition (m4)
|-- training/
|   |-- train_qwen35_lora.py               LoRA / partial-FT trainer (TRL >=0.18)
|   |-- eval_qwen35_lora.py                holdout evaluator (--adapter optional)
|   |-- patch_qwen35_template.py           chat-template generation-marker patcher
|   |-- chat_with_adapter.py               interactive + probe-set generality tester
|   |-- rejudge_failed_episodes.py         post-hoc rejudge for transient judge errors
|   `-- compare_pre_post.py                per-skill diff utility
|-- docker/                                vast.ai training image + setup docs
|-- data/pipeline-runs/default/            full run artifacts
`-- b1.reports/                            engineering log (yymmdd.title.txt)
\end{verbatim}

\section{Per-skill v2.0 result ($n{=}5$ per skill, sorted by $\Delta$)}\label{app:per-skill}
Lift cluster (procedure helped):
\begin{itemize}[leftmargin=*,nosep]
  \item accident-fallacy $+0.40$, spatial-reasoning $+0.40$
  \item false-dilemma, confirmation-bias, anchoring-bias, appeal-to-emotion, base-rate-neglect, transitive-inference, necessary-and-sufficient-conditions, metaphor, emotional-self-regulation: $+0.20$ each
\end{itemize}
Flat cluster ($\Delta = 0$): 25 skills, of which 22 are at \BL{} $\geq 0.80$ (ceiling-bound).

Regression cluster (procedure hurt; all $5/5$ \BL{} $\to 4/5$ \CU): complex-question, hypothetical-syllogism, framing-effect, equivocation.

The regression cluster overlaps with deductive-logic skills where the base 4B has ceiling competence and the 5-step procedure introduces sub-inference error opportunities --- the same lift/flat/regression pattern observed at 0.8B in v1, shifted upward by base scaling.

\section{Engineering reports cited}\label{app:reports}
Full text of these reports is in the repository's \texttt{b1.reports/} directory:
\begin{itemize}[leftmargin=*,nosep]
  \item \texttt{260426.findings-pre-post-sft-iterations.txt} --- v1--v1.8 consolidation
  \item \texttt{260428.sft-v1\_9-2b-result.txt} --- v1.9 result + attribution
  \item \texttt{260428.v1\_9-generality-probe.txt} --- v1.9 generality regression
  \item \texttt{260428.sft-v2\_0-4b-result.txt} --- v2.0 result + attribution caveat
  \item \texttt{260428.v2\_0-generality-probe.txt} --- v2.0 generality preservation
  \item \texttt{260506.bench-format-sensitivity-finding.txt} --- pre-SFT 4B format artifact
\end{itemize}

\end{document}
